\section{太极（Taichi）：一个人的编程语言}

\subsection{什么是计算机图形学}

胡渊鸣对计算机图形学的定义简洁有力：\textbf{用加减乘除创造一个虚拟世界}。CPU 和 GPU 都只能做加减乘除和三角函数，但却要用这些最基本的 building blocks 创建栩栩如生的游戏画面。

图形学的发展大致分几个阶段：

\begin{enumerate}
\item \textbf{可视化与真实感渲染}（约1995年前）
\item \textbf{物理仿真}（1995年后）：Fedkiw 等人的烟雾模拟（Visual Simulation of Smoke）开启了仿真时代
\item \textbf{计算摄影学}：super resolution 等技术，与 CV 不分家
\item \textbf{Neural Graphics}：DLSS、NeRF、Gaussian Splatting 等 AI 驱动的新范式
\end{enumerate}

\begin{knowledgebox}{判断一个领域是否在衰退的信号}
"当你发现一个问题越来越干净、越来越被 well-defined 的时候，就说明这个问题没什么可做的了。当它形成了特别系统的 benchmark，当解法搞得越来越复杂的时候，说明开疆拓土的机会已经比较少了。"——胡渊鸣以渲染领域的 AMCMCPPM（自适应马尔科夫链蒙特卡罗统一路径采样）为例说明这一点。
\end{knowledgebox}

\subsection{太极的诞生与使命}

太极的核心使命是构建一套\textbf{基础设施（Infrastructure）}，使得仿真能够更容易开发、运行更快、更省内存，并与周围生态（Python、PyTorch）友好集成。

一个标志性成就是在单块 3090 GPU（24GB 显存）上运行10亿（$10^9$）粒子的仿真。为实现这一点，每个粒子占用的存储空间不能超过24字节。他在编译器层面做了极端的量化优化——一个粒子的 XYZ 坐标用一个32位整数表示（X 用11位，Y 用11位，Z 用10位），这就是 QuantTaichi 的工作。

\begin{importantbox}{太极为什么需要自建编译器}
关键在于 CUDA 虽然让 C++ 程序员能编写 GPU 程序，但 Python 生态、LLVM 编译器基础设施、MPM 等仿真方法的成熟，创造了在更高层次构建领域专用语言的机会。太极基于 Python 语法，通过自建编译器将高级代码编译为高性能 GPU 程序，用户无需手动处理位操作和内存优化。
\end{importantbox}

在 MIT 独立开发太极的体验被胡渊鸣描述为"很爽"——他同时扮演了四个角色：

\begin{table}[H]
\centering
\caption{太极开发中的四重角色}
\begin{tabular}{ll}
\toprule
\textbf{角色} & \textbf{具体工作} \\
\midrule
产品经理 & 分析仿真用户需求（因为自己就是用户） \\
科学家 & 了解前沿技术，将其转化为产品特性 \\
工程师 & CI/CD、软件架构、C++性能优化、多线程编译、编译缓存 \\
CEO & 整合三个角色的成果，让更多人用上（知乎宣传等） \\
\bottomrule
\end{tabular}
\end{table}

"用自己的编程语言去实现自己的程序，然后发现还比 CUDA 好用的时候，我觉得那时候有点牛逼坏了。"——这是太极开发过程中最让他激动的时刻。

\subsection{99行代码实现冰雪奇缘}

太极最广为人知的 demo 是"99行代码实现冰雪奇缘效果"——用极短的 Python 代码实现了材料点法（MPM）驱动的雪花仿真。这个 demo 在社交媒体上获得了巨大关注，也成为了胡渊鸣的标志性标签之一。

然而他对这个标签的态度是复杂的："我相信你已经对这种标签早就去魅了，并且这样的标签完全不能展示真实的你。"事实上，99行 demo 只是太极能力的冰山一角——它背后是完整的编译器栈、类型系统、自动微分引擎、多后端代码生成器等大量基础设施。

\begin{knowledgebox}{太极的可微编程（Differentiable Programming）}
太极的一个关键特性是支持可微编程——即你的仿真程序可以自动求导。这意味着你可以：
\begin{itemize}
\item 给定一个目标状态，反向优化仿真参数
\item 将 simulator 嵌入到 neural network 的训练循环中
\item 用 gradient descent 来优化 robot 的控制策略
\end{itemize}
DiffTaichi 和 ChainQueen 两篇论文正是基于这个特性，将 simulation 和 robotics 结合起来。这个方向是清华学长家俊建议的——"你要不试试把你的 simulator 搞成 differentiable？"
\end{knowledgebox}

胡渊鸣回忆说，只需在太极代码中加上一两行代码，就能把普通仿真变成可微仿真——"我也觉得只有自己造了一个编译器才能做到这个事情。"这种编译器级别的抽象能力，是手写 CUDA 代码无法企及的。

\subsection{太极 2.0 的愿景}

访谈中透露，太极 2.0 的规划正在进行中。目标是让太极在大模型时代仍有广泛应用，影响力达到1.0的50倍以上。具体方向可能包括：

\begin{itemize}
\item \textbf{支持 Transformer 优化}：如果当初太极更早向这个方向发展，"可能现在太极是 Triton 的位置"
\item \textbf{Hybrid Simulator}：将 simulation 和 neural network 真正结合，既有物理仿真部分又有 data-driven 部分
\item \textbf{LLM 生态支持}：在 LLM 领域找到太极的应用场景
\end{itemize}

一个有趣的数据：太极的 GitHub Stars 数目超过了 NVIDIA Warp、Triton、Mojo 等项目——"太极还是很 popular 的，只不过没有找到一个大的领域去施展。"

\subsection{开源哲学的三次演变}

胡渊鸣对开源的看法经历了三个阶段：

\textbf{第一阶段（MIT 博士期间）}：开源大法好。当时很多 Graphics 论文没有开源代码，他坚持让太极的所有实验结果都可通过一行命令复现。

\textbf{第二阶段（创业早期）}：开源商业化九死一生。"你都把最好的东西给人家了，那怎么赚钱呢？"太极投入减少，团队调去做 Meshy。

\textbf{第三阶段（当前）}：开源是战略性的 marketing 工具。

\begin{importantbox}{开源作为技术营销}
"开源确实不挣钱，但一个好的开源项目能帮你吸引到最好的人才和最好的客户——它变成了一个 marketing 手段。"在游戏行业，研发和市场推广投入往往是1:1甚至更多。通过开源做技术营销，对于不擅长传统营销的技术创始人来说是最佳路径。但这需要"非常 strategic"——必须想到一个好的商业化方式与之互补。
\end{importantbox}

他坦言，开源是他作为 Founder 理想中不可甩掉的部分——"迟早你可能又会走到开源这条路上"。目前团队正在规划太极 2.0，定位为 "infrastructure supporting Graphics and AI"。

\subsection{本章小结}

太极项目体现了胡渊鸣"从底层开始造轮子"的极客精神，也是他从学术走向创业的桥梁。技术上它展示了领域专用编译器的巨大潜力（10亿粒子仿真），方法论上它训练了"产品经理+科学家+工程师+CEO"的全栈能力。开源哲学的三次演变则反映了理想与商业现实之间的平衡。

